\chapter{Penutup}

\section*{Sintesis Isi Buku}\addcontentsline{toc}{section}{Sintesis Isi Buku}

Empat belas bab sebelumnya menempuh satu jalur, dari sistem yang dijalankan
menuju model yang membangkitkan sistem itu kembali. Bab 1 sampai Bab 3
memperlakukan sistem sebagai objek yang dapat dijalankan dan diamati, yaitu
pipeline pemrosesan video yang dieksekusi dengan kontainer, diukur melalui metrik
dan dasbor, lalu dibatasi perilakunya dengan kontrol keamanan pada saat berjalan.

Bab 4 sampai Bab 6 memindahkan pekerjaan manual itu ke otomasi. Pipeline CI/CD
menjalankan pengujian dan penerapan tanpa langkah manual, Terraform menyatakan
infrastruktur sebagai kode deklaratif untuk tiga lingkungan, dan Kubernetes
menjalankan sistem yang sama di atas orkestrator. Pada titik ini jumlah artefak
operasional meningkat, dan duplikasi antar lingkungan menjadi masalah yang nyata.

Bab 7 sampai Bab 13 menjawab masalah tersebut dengan menaikkan tingkat
abstraksi. Model lingkungan dan templat membangkitkan berkas Terraform dan
Kubernetes, metamodel MediaFlow memformalkan domain pipeline media, DSL tekstual
dan visual menyediakan dua sintaks konkret di atas sintaks abstrak yang sama,
transformasi model ke model memetakan model ke representasi lain, generator model
ke teks menurunkan artefak yang dipakai langsung oleh penampil, dan aturan
validasi menolak model yang salah sebelum artefaknya dibangkitkan. Bab 14
menutup rangkaian itu dengan sistem multi-agen yang menjalankan perkakas tersebut
melalui antrean tugas bersama.

\section*{Jawaban atas Rumusan Masalah}\addcontentsline{toc}{section}{Jawaban atas Rumusan Masalah}

Pendahuluan mengajukan tiga pertanyaan. Berikut jawaban yang diberikan isi buku
ini.

\begin{enumerate}
  \item \textbf{Jalur dari sistem berjalan menuju model yang dapat dieksekusi.}
        Setiap abstraksi diperkenalkan setelah masalah yang diatasinya terlihat
        pada sistem yang berjalan. Metamodel diperkenalkan setelah pembaca
        menulis berkas konfigurasi berulang, dan generator diperkenalkan setelah
        pembaca menyalin artefak yang sama ke tiga lingkungan. Urutan ini
        menjaga kaitan antara model dan artefak operasional, karena keduanya
        dibandingkan pada sistem yang sama.
  \item \textbf{Satu studi kasus lintas tahapan.} Domain MediaFlow dipakai
        berturut-turut pada Bab 8 sampai Bab 14, sedangkan \textit{stack} layanan
        text to speech dipakai pada Bab 4 sampai Bab 7. Dengan cara itu,
        perbedaan antar pendekatan dapat dinilai tanpa mengganti domain, dan
        akibat satu keputusan perancangan metamodel dapat ditelusuri sampai ke
        artefak yang dibangkitkan.
  \item \textbf{Pengurangan duplikasi dan batasnya.} Pendekatan berbasis model
        memangkas duplikasi ketika artefak yang sama perlu diturunkan untuk
        banyak lingkungan atau banyak target, sebagaimana ditunjukkan pada Bab 7
        dan Bab 12. Biayanya muncul pada sisi lain, yaitu metamodel, templat, dan
        generator yang harus dipelihara, serta kesalahan generator yang langsung
        menyebar ke seluruh keluaran. Untuk sistem kecil dengan satu lingkungan,
        menulis artefak secara langsung masih lebih murah.
\end{enumerate}

\section*{Kontribusi}\addcontentsline{toc}{section}{Kontribusi}

Buku ini memberikan tiga kontribusi. Pertama, satu rangkaian materi yang
menggabungkan praktik DevOps \cite{bass2015devops, humble2010cd, morris2021iac}
dengan rekayasa berbasis model \cite{schmidt2006mde, brambilla2017mdse} pada satu
jalur yang berangkat dari eksekusi. Kedua, satu metamodel beserta perkakas
pendukungnya, yaitu DSL tekstual, DSL visual, transformasi ETL
\cite{kolovos2008etl}, generator EGL \cite{rose2008egl}, dan aturan validasi EVL
\cite{kolovos2009evl}, yang seluruhnya dapat dijalankan ulang dari repositori
publik. Ketiga, perbandingan antar perkakas pada tugas yang sama, sehingga
pembaca memiliki dasar untuk memilih perkakas pada proyeknya sendiri.

\section*{Keterbatasan}\addcontentsline{toc}{section}{Keterbatasan}

Empat keterbatasan perlu dinyatakan terbuka.

\begin{enumerate}
  \item \textbf{Belum ada pengujian empiris pada pembelajaran.} Pengaruh urutan
        materi ini terhadap capaian belajar mahasiswa belum diukur dengan
        rancangan penelitian yang terkendali.
  \item \textbf{Skala studi kasus.} Seluruh studi kasus dijalankan pada satu
        mesin dengan beban kecil. Perilaku pada klaster produksi, beban nyata,
        dan model berukuran besar belum diuji.
  \item \textbf{Keterikatan pada versi perkakas.} Perintah dan keluaran
        mengikuti versi perkakas yang dipakai saat penulisan. Perubahan versi
        Terraform, Kubernetes, Eclipse, atau Epsilon dapat mengubah rincian
        langkah, sehingga contoh perlu dijalankan ulang pada cetakan berikutnya.
  \item \textbf{Perbandingan yang bersifat kualitatif.} Perbandingan antar
        perkakas disusun dari pengalaman menerapkannya pada satu tugas yang
        sama, dan belum memakai pengukuran berulang dengan analisis statistik.
\end{enumerate}

\section*{Arah Pengembangan}\addcontentsline{toc}{section}{Arah Pengembangan}

Keterbatasan di atas menunjukkan pekerjaan lanjutan yang dapat dikerjakan.

\begin{enumerate}
  \item \textbf{Pengujian di kelas.} Urutan materi ini dapat diuji dengan
        membandingkan dua kelompok kelas, memakai nilai sebelum dan sesudah,
        kualitas artefak yang dihasilkan, serta waktu yang dibutuhkan mahasiswa
        untuk menyelesaikan tugas yang sama.
  \item \textbf{Pengukuran kuantitatif antar perkakas.} Perbandingan pada Bab 9,
        Bab 11, dan Bab 13 dapat diulang dengan model berukuran lebih besar dan
        pengulangan yang cukup, sehingga waktu eksekusi, jumlah baris, dan
        tingkat kesalahan dapat dilaporkan dengan ukuran efek.
  \item \textbf{Perluasan target generasi.} Generator pada Bab 12 dapat
        diperluas untuk membangkitkan manifes Kubernetes dan berkas Terraform
        secara langsung dari metamodel MediaFlow, sehingga rantai dari model ke
        sistem berjalan tertutup penuh.
  \item \textbf{Peran model bahasa besar.} Bab 14 membuka kemungkinan agen
        menjalankan rantai perkakas ini. Langkah lanjutannya adalah mengukur
        keandalan agen pada tugas pemodelan, termasuk tingkat keberhasilan,
        jumlah perbaikan yang dibutuhkan, dan biayanya.
\end{enumerate}

\section*{Ringkasan}\addcontentsline{toc}{section}{Ringkasan}

Buku ini berangkat dari satu gagasan sederhana: sistem perangkat lunak lebih
mudah dipahami ketika dijalankan lebih dahulu, dan abstraksi lebih mudah diterima
ketika masalah yang diatasinya sudah terlihat. Dari gagasan itu disusun jalur
yang bergerak dari eksekusi, observasi, dan otomasi, menuju model, transformasi,
validasi, dan generasi. Seluruh artefaknya tersedia terbuka, sehingga pembaca
dapat menjalankannya kembali, memeriksa klaim yang dibuat, dan memakainya sebagai
titik awal pada pekerjaannya sendiri.
